\chapter{Preface}

Most courses on computer architecture are built from the bottom up. They begin
with the bit, climb through number systems, Boolean algebra and logic networks,
assemble a processor out of gates, and only in the last weeks let the student
write a program for the machine that the previous months have described. The
order is logically impeccable, and for many first-year students it is also the
reason why the subject feels remote: for ten weeks every answer arrives before
its question.

This book is the written form of a course that turns the order round without
giving up any of the rigour. \emph{Program first, formalise later}: in the
first week the reader installs a small console, types a dozen lines of \riscv
assembly and sees a rectangle move across a screen. From then on every formal
topic is introduced by something that the reader's own program has already
done. Two's complement arrives after a position has walked off the left edge
of the screen and the debugger has shown \hex{FFFFFFFF}. Boolean algebra
arrives after a mask has been used to ask whether a button is pressed. The
flip-flop arrives after a counter has survived from one frame to the next, and
the instruction cycle after the reader has pressed \key{Step} a few hundred
times and watched the program counter advance by four.

Nothing is removed from the classical syllabus. The reader will find
positional notation and the IEEE~754 formats, axioms and theorems of Boolean
algebra with their proofs, canonical forms and Karnaugh maps, the transistor as
a switch, adders, multiplexers and decoders, latches and flip-flops, finite
state machines, registers and buses, memory hierarchies and permanent storage,
the datapath and the control unit, interrupts, the video card and the
peripheral devices. What changes is the direction of the argument: each
chapter opens on the machine, leaves it for as long as the theory requires, and
comes back to it with a section called \emph{Back to the machine}, where the
same rules are read again in the instructions of a real processor.

\section*{The machine}

The processor is a \riscv core: the RV32IMAC instruction set, as implemented by
the Espressif ESP32\nobreakdash-C6. \riscv is an open standard, small enough
to be learnt completely in one semester and real enough to be found in
billions of devices; it replaces here the Motorola~68000 of the previous
editions of the course, which survives only in a few historical comparisons.

The environment is \prg, an open educational platform developed at the
University of Naples ``Parthenope''. A \prg program is a \term{cartridge}: a
small file of native machine code that supplies three functions, one called
once and two called at every frame. The same cartridge runs on a breadboard
kit, on the QEMU emulator and on \prgqt, a desktop application whose built-in
debugger shows the thirty-two registers and the memory while the program runs
one instruction at a time. No board is needed to follow the book: a laptop is
enough.

\section*{The project}

A platform that is used only for exercises is a gymnasium. \prg is also a
workshop, and the book has a project that runs through it: \vermone, read
\emph{vermone}, ``the big worm'', a tribute to Snake written from an empty file
in \riscv assembly. After every block of four chapters a tutorial chapter adds
a part of the game, using only what those four chapters have taught: first a
head that crosses a world without edges, then steering, food and a score, then
a body that grows and the rules that end a game, then digits, levels and
sound. A last tutorial packages the finished game for the \prg Cartridge
Store. Every tutorial prints every line of the program, tells which lines are
new and which have changed, and stops at checkpoints where the program must
run and behave in a stated way.

\section*{What ``verified'' means in this book}

A listing that does not assemble teaches the wrong lesson. The assembly code
in this book is not typed into the text: it is read from files that are
assembled, run on the upstream QEMU and on the headless runner of \prgqt, and
compared with expected results by automatic tests; the four parts of \vermone
are played by a scripted player that also takes the screenshots printed in the
tutorial chapters. Where a fragment could not be run on the author's machine
(the CUDA and Verilog fragments of the last chapter, the installation steps for
operating systems other than the one used to write the book) the text says so.

\section*{Acknowledgements}

Mariacarla Staffa teaches the laboratory half of the course, where the
exercises of this book are done with pencil, paper and keyboard. Ivan Cafiero
and Simone Boscaglia contributed to \prg as students and showed that the
platform can be understood, and extended, by the people for whom it is meant.
Giuseppe Coviello contributed the reference implementation of \vermone: the
four programs that the tutorial chapters build line by line, and that the
automatic tests play, are his.
The students of the previous editions of the course wrote, with their
questions, most of the paragraphs that begin with ``Watch out''.

\begin{flushright}
Raffaele Montella\\
Naples, October 2026
\end{flushright}
